% ARGUES: what the four experiments jointly license, and no more.
% DRAWS FROM: results, once they exist.
% GATED ON: the confirmatory fit.

\section{Discussion}
\label{sec:discussion}

% REFRAMED 2026-08-31 after the confirmatory fit. The chain below previously led
% with E1 and asserted a differential cascade. E1's registered interaction did
% not reject. Do not restore the old ordering, and do not restate E1's main
% effect as though it were the hypothesis that was tested.

\guide{The thesis chain, reordered to match what the experiments returned, and
stated as a chain so a reader can see which link is weakest.

E2 carries the paper. A prompt cue moves fabrication within a fixed policy, and
it raises fabrication in post-trained checkpoints far more than in base ones.
On the original twelve it drives the base models down; the fifteen-model refit
adds a base model that goes up, so lead with the post-trained rise and the
interaction, not with the base fall. That is the strongest and cleanest result
and it should lead.

E1 is a large main effect whose registered interaction did not reject. Post-
training raises fabrication enormously, and it does not raise it preferentially
in the cell where capability is held roughly constant, which is what H1.1
predicted. Report the null as a null. The honest sentence is that post-training
raises fabrication across the board rather than selectively where fabrication is
most tempting, and the co-primary intersection-union rule therefore returns a
loss for the headline T1 claim.

E3 asks whether the signal is internally represented. E4 asks whether a
plausible surface proxy prefers the higher-fabricating policy, and its answer is
qualified by the count-matched control: the preference is largely ``cites
more.''

One interpretive paragraph, clearly labeled as interpretation rather than
result. A broad rise rather than a targeted one is what
\citet{qi2026rewardseeker} report for reward-seeking induced deliberately: the
disposition generalized well outside the environments that produced it instead
of firing only where it had been trained. Our E1 null is consistent with a
general disposition rather than a strategy that activates on detecting a hard
question. This is a reading of a null result against an external parallel and
carries no evidential weight; write it as such or cut it.

The negative stage-by-unanswerable interaction wants the same treatment. A
ceiling reading is available, since fabrication on unanswerable items starts
high and has less room to climb, but the design does not test it.

Keep the boundary paragraph, and keep it early. This design does not identify
the reward signal that produced the behavior. It shows that fabrication moves
with post-training stage and with a prompt cue, and that a surface proxy would
reward it. Inferring that RLHF specifically caused it requires an assumption the
design does not test, and the E1 null makes the tempting causal story weaker
rather than stronger. A reviewer will look for that overreach first.

Write the boundary paragraph early rather than at the end. This design does not
identify the reward signal that produced the behavior. It shows that fabrication
moves with post-training stage and with a prompt cue, and that a surface proxy
would reward it. Inferring that RLHF specifically caused it requires an
assumption the design does not test. A reviewer will look for that
overreach first.

If the dose-response convergence with the concurrent solution-hacking work
holds, it is worth a short paragraph: two designs, different constructs, same
gradient with difficulty.

Practical implication, kept to what follows from the data: evaluation and
training setups that reward citation presence without checking citation
existence select for exactly this behavior, and the E4 proxy is a direct
demonstration if it lands.}

\paragraph{What this design cannot do, first.}
This study does not identify the reward signal that produced the behavior. It
shows that citation fabrication moves with post-training stage, that a prompt cue
raises it in post-trained checkpoints far more than in base ones, and that a
naive surface proxy would reward the post-trained policy.
Concluding that RLHF specifically caused any of it requires an assumption the
design does not test: we observe checkpoints along a training pipeline, not an
intervention on the reward. The E1 null makes that causal story weaker rather
than stronger, because the effect we can measure is a broad rise rather than the
targeted one the reward-hacking account predicted. We state this at the top
rather than at the end because every paragraph below inherits it.

\paragraph{E2 carries the paper.}
The clearest result is that a prompt cue moves fabrication within a fixed policy,
and moves the two populations in opposite directions. Asking for citations drives
base checkpoints \emph{down} (9.04\% to 0.77\%) and post-trained checkpoints
\emph{up} (34.18\% to 73.70\%), a 47.8-point interaction that holds in all three
domains and survives under both analysis pipelines. Because the weights are fixed
while the prompt changes, this is not a claim about what training did; it is a
demonstration that the disposition is elicitable, and that whatever post-training
installed, it responds to being asked for sources by producing more of them
whether or not they exist.

The fifteen-checkpoint refit (Section~\ref{sec:refit15}) keeps the interaction,
41.5 points with Holm-adjusted $p = 3 \times 10^{-31}$, and keeps the post-trained
rise almost unchanged, 35.16\% to 74.91\%. What it does not keep is the base-side
fall as a general statement. Three of the four base checkpoints fabricate less
when asked for citations, but the Qwen2.5-72B base goes from 47.8\% to 71.7\%,
responding to the cue the way post-trained checkpoints do. A model released as a
base checkpoint is not guaranteed to behave like one, and we do not know what in
that checkpoint's training produces this. What holds on both rosters is the
post-trained rise and the interaction, and that is the claim we make. One
qualification belongs beside it. Adding citation count as a covariate shrinks
the fifteen-model interaction from an odds ratio of 6.75 to 2.28, so a good part
of the cue effect runs through post-trained checkpoints citing more, the same
quantity E4 finds a naive proxy rewarding. That adjusted figure is not a direct
effect. Citation count is itself moved by the cue, so the covariate sits on the
causal path, and the outcome counts a generation as fabricating if any one of its
citations is fabricated, so more citations raise it mechanically. The adjusted
odds ratio says how much of the interaction survives holding count fixed; it does
not split the effect into a count part and a propensity part. The interaction
stays above one after adjustment on both pipelines (2.28 corrected, 1.48
registered).

The de-confounding control sharpens this. A reasonable objection is that
\texttt{citations\_expected} simply asks for a longer answer, and longer answers
afford more citation slots. The length-matched control answers it: the
citation-specific cue exceeds the pure length demand by 28.35 points on the
corrected pipeline and 12.64 on the registered one, both far inside the
pre-committed ``citation cue drives it'' bin. That rule was fixed before any
Phase C number existed and was written to accommodate the \emph{opposite} result,
since the exploratory pilot had pointed the other way. It did not fire.

\paragraph{E1 is a large main effect whose registered interaction did not reject.}
Post-training raises fabrication enormously: 4.78\% at base to 56.66\% at the
final stage, an odds ratio of 2.41 per stage. But H1.1 predicted something more
specific, that the rise would concentrate in the obscure-but-real cell where
capability is roughly held constant and fabrication is most tempting. It did not.
The registered interaction is 0.084 with a Holm-adjusted $p$ of 0.360, and the
co-primary binary contrast does not reject either, so the intersection-union rule
returns a loss for the headline T1 claim under both pipelines. We report that as a
null, not as a trend, and not as a smaller version of the predicted effect.

\textbf{That null is uninformative rather than informative, and the reason is a
design shortfall we have to declare.} The pre-registration's power figures for E1
(97.5\% at $\alpha = 0.05$, 90.0\% at the Holm floor) were computed on a
23-checkpoint roster across six families. The study ran twelve checkpoints across
three, having not collected the 405B anchor or the cross-family breadth arm.
Re-running the registered power machinery on the roster that was actually fitted,
with the effect pinned to the same pre-registered declared value so the
calculation stays a-priori rather than a restatement of the $p$-value, gives
\textbf{73.0\% $[66.28, 79.02]$ at $\alpha = 0.05$ and 54.0\% at the Holm floor}.
Both fail the pre-registered 80\% gate. So the design as run would miss an effect
of the registered size roughly a quarter of the time at nominal $\alpha$, and
close to half the time at the threshold this interaction actually faces as the
second member of its Holm family. The same recomputation for E2 gives 73.0\% and
51.0\%, against registered figures of 95.5\% and 85.0\%. That bears on none of
the E2 conclusions, because power governs how a failure to reject should be read
and E2 rejects at Holm $p = 2 \times 10^{-34}$; we report it because the
registered figure describes a roster we did not run.

The correct reading is therefore not that post-training fails to raise fabrication
selectively, but that \textbf{this experiment was not powered to settle the
question at the effect size it registered}. We flag two things that reading must
not be stretched into. It does not rescue T1: low power explains how a null can
arise, it is not evidence for the effect, and the registered conclusion remains
``not supported.'' And it is not a claim about power at the observed estimate,
which would be a circular restatement of the same $p$-value. What we can say
positively is the main effect, which is large and unambiguous, and which no
amount of power discussion changes: post-training raises fabrication across the
board.

The fifteen-checkpoint refit moves this reading. Collecting three registered
checkpoints under the pre-committed rule of amendment \#13 lifts a-priori power
for the E1 interaction to 83.0\% at nominal $\alpha$, above the registered gate,
though still not at the Holm floor (69.0\%). On that roster the interaction is
null again, Holm-adjusted $p = 0.192$, and its point estimate turns slightly
negative ($\mathrm{OR} = 0.93$). Its 95\% interval, $[0.86, 1.01]$ (registered
$[0.85, 1.02]$), is the more useful number: whatever positive interaction the
data still allow is at most one or two percent in odds per stage. The
stage-as-factor sensitivity puts every stage-by-obscure-real term below one. Amendment \#13 fixed in advance that a null
at adequate power would be reported as evidence against selective concentration,
and we do so with the qualifier the Holm floor requires: at the nominal level the
fifteen-model design was powered for the registered effect and did not find it,
and the direction it did find is the opposite one. The binary base-versus-post
contrast does reject on fifteen checkpoints, but the intersection-union rule
needs both components, so T1 remains unsupported.

Two further features of E1 are worth naming and not over-reading. Treating stage
as a factor shows most of the movement is already present at SFT, with DPO and
RLHF indistinguishable from each other, so the cascade is not a story about the
final preference-optimization step. And the stage-by-unanswerable interaction is
negative. A ceiling reading is available, since fabrication on unanswerable items
starts high and has less room to climb, but the design does not test it and we do
not assert it.

\paragraph{An interpretation, labeled as one.}
A broad rise rather than a targeted one is what \citet{qi2026rewardseeker} report
for reward-seeking induced deliberately: the disposition generalized well outside
the environments that produced it, instead of firing only where it had been
trained. Our E1 null is consistent with a general disposition rather than a
strategy that activates on detecting a hard question. This is a reading of a null
result against an external parallel. It carries no evidential weight and nothing
later depends on it.

\paragraph{E4 shows the proxy would reward it, and shows exactly which feature.}
A preference model trained only on surface citation features prefers post-trained
output 85.8\% of the time on cross-family transfer, which is the plausibility
argument the thesis needs: a reward signal that keys on citation presence would
select for the higher-fabricating policy. The count-matched diagnostic tells us
which feature carries it. Holding citation count exactly equal collapses the
preference to 51.6\%, while the placebo lands at 50.3\% with an interval
straddling chance. So the proxy is reading \emph{quantity}, not plausibility of
apparatus. We therefore claim that citation count is rewarded, and we do not
claim that a naive judge is fooled by fabricated-looking scholarly detail.
Quantity is an entirely ordinary thing for a reward model to key on, which is
what makes the mechanism plausible rather than exotic.

\paragraph{What the numbers will and will not support.}
The base-versus-post-trained \emph{contrast} is robust: it holds under both
pipelines, in every domain, and under the UNVERIFIABLE-as-fabricated recoding
that excludes no generation at all. The \emph{absolute} rates are not. They rest
on an LLM judge measured at 80\% precision and 0.781 recall on FABRICATED, so a
reader should treat ``56.66\% of final-stage generations contain a fabricated
citation'' as an estimate with real measurement error, and treat the direction
and size of the differences as the durable finding. All human labels are the
author's, so there is no inter-annotator reliability and no Cohen's $\kappa$.
Nothing here is proven; a pre-registered study with one instrument and one
annotator does not prove things.

Two scope limits deserve stating in the same breath, because both cut against
the tempting summary. The closed-API block is empty, so \textbf{no claim about
frontier closed models follows from this study} and H1.2b is untested rather than
unsupported. And base models are not the honest comparators a casual reading
imagines: the three original base checkpoints cite \emph{less}, not more
accurately, and under the citation cue they mostly stop citing altogether. Their
column is a floor produced partly by non-compliance with the instruction, not a
demonstration of restraint. The fourth base checkpoint, Qwen2.5-72B, cites
heavily and fabricates at post-trained rates, which is a second reason not to
treat ``base'' as a behavioral category.

\paragraph{Practical implication.}
The narrow version, which is all the data supports, is this: an evaluation or
training setup that rewards citation presence without checking citation existence
selects for exactly the behavior measured here. E4 is a direct demonstration that
a plausible proxy does so, and it identifies count as the operative feature.
Checking existence is cheap relative to the alternative, and the failure it
prevents is the one that survives review.

\subsection{Future work}

Four things this study sets up and does not deliver.

\textbf{A second annotator.} Every human label here is the author's, so the
obvious next step is an independent annotator on a subset and a Cohen's $\kappa$.
We note one design constraint learned the hard way: the annotator has to be able
to judge whether a Federal Reporter citation resolves to a real case. An
annotator who cannot would produce a low $\kappa$ measuring their unfamiliarity
rather than the protocol's ambiguity, which is worse than reporting no $\kappa$
at all, because it reads as evidence that the dividing rules are unclear when
they are not.

\textbf{A design powered for the interaction.} The registered roster was reduced
before collection, and the realized design reached 73\% power for the E1 effect
at nominal $\alpha$ and 54\% at the Holm floor; the fifteen-checkpoint refit
reached 83\% and 69\%. Settling the question at the Holm floor needs the breadth
the registration called for: more families carrying both base and post-trained
checkpoints, and enough of them that the family-level random effect is estimated
from more than four levels.

\textbf{Retrieval-augmented settings.} It is tempting to assume retrieval
dissolves this problem. \citet{magesh2024} report that retrieval-based legal
research tools still hallucinate between 17\% and 33\% of the time, so the
assumption needs testing rather than granting. The design here transfers
directly: the same question set, the same conditions, with retrieval in the loop.

\textbf{Whether the signal is internally represented, and whether that is
usable.} E3 (Section~\ref{sec:e3}) ran on the 8B scale-check arm only. There, a
fabrication-correlated signal is linearly decodable at the answer tokens from
disjoint questions, and it survives the registered difficulty deconfound; the
general-truthfulness probe, the registered primary, does not transfer at all.
Three things remain open. The 70B primary is uncollected, so whether the
within-domain result holds at the registered scale is unknown. The activations
were read from 4-bit weights. And decodability is detection, not intervention:
whether the signal can be steered, rather than merely read, is the more useful
question and this design does not ask it.
